\section{Theoretical Setup}
\label{sec:setup}

\subsection{Problem setting}

Let $n\in\mathbb N_{\geq1}$, let
$\mathbb N_0=\{0,1,2,\ldots\}$, and put
$\mathcal X=\{-1,+1\}^n$ and
$\mathcal H\subseteq\{-1,+1\}^{\mathcal X}$.  We write
$\Delta(\mathcal X)$ for the probability distributions on $\mathcal X$ and
$[m]=\{1,\ldots,m\}$.  Fix the source tie label
$s_0=\operatorname{sign}(0)\in\{-1,+1\}$ and define
\[
\operatorname{sign}_{s_0}(z)=
\begin{cases}
+1,&z>0,\\
-1,&z<0,\\
s_0,&z=0.
\end{cases}
\]
For a map $\varphi:\mathcal X\to\mathbb R^d$, a coefficient
$w\in\mathbb R^d$, $\mathcal D\in\Delta(\mathcal X)$, and
$h\in\mathcal H$, define the tie-resolved linear error
\[
R_{\mathcal D,h}(w,\varphi)
:=\Pr_{x\sim\mathcal D}\!\left[
\operatorname{sign}_{s_0}(\langle w,\varphi(x)\rangle)h(x)<0
\right].
\]
For $\alpha\geq0$, the probabilistic dimension complexity is
\[
\operatorname{dc}_{\alpha}(\mathcal H)
:=\min\left\{d\in\mathbb N_0:
\begin{array}{l}
\exists\text{ a law }\mathcal P\text{ over }
\varphi:\mathcal X\to\mathbb R^d\quad
\forall\mathcal D\in\Delta(\mathcal X)\quad\forall h\in\mathcal H,\\[1mm]
\displaystyle
\mathbb E_{\varphi\sim\mathcal P}\!\left[
\inf_{w\in\mathbb R^d}R_{\mathcal D,h}(w,\varphi)
\right]\leq\alpha
\end{array}
\right\}.
\tag{1}
\]
Thus the feature-map law is selected before every later distribution and
target, whereas the optimizing coefficient may depend on the realized map,
distribution, and target.  The fixed tie label makes the definition
nonvacuous at a zero score.

Fix a depth $L\in\mathbb N_{\geq1}$, positive integer widths
$n_0=n,n_1,\ldots,n_{L-1},n_L=1$, and matrices
$\theta_\ell\in\mathbb R^{n_\ell\times n_{\ell-1}}$.  The network has no
biases and has parameter count
\[
S:=\sum_{\ell=1}^{L}n_\ell n_{\ell-1},
\qquad
\theta=(\theta_1,\ldots,\theta_L)\in\mathbb R^S.
\tag{2}
\]
With $\sigma(a)=\max\{0,a\}$ coordinatewise, define
\[
z_0(\theta,x)=x,
\qquad
u_\ell(\theta,x)=\theta_\ell z_{\ell-1}(\theta,x),
\qquad
z_\ell(\theta,x)=\sigma(u_\ell(\theta,x))
\quad(1\leq\ell<L),
\]
and $f_\theta(x)=\theta_Lz_{L-1}(\theta,x)$.  The coordinate
$u_{\ell,j}$ denotes the $j$th entry of $u_\ell$.

Fix the source protocol's ReLU-kink convention $\kappa\in[0,1]$ and write
\[
\rho_\kappa(a)=
\begin{cases}
1,&a>0,\\
0,&a<0,\\
\kappa,&a=0.
\end{cases}
\]
The notation $\nabla_\theta^{(\kappa)}$ denotes the resulting single-valued
back-propagation derivative.  Both $s_0$ and $\kappa$ are fixed protocol
conventions, not choices made by the theorem.

Initialize all weights independently according to
\[
(\theta_\ell^{(0)})_{jk}\sim\mathcal N(0,1/n_{\ell-1}),
\qquad 1\leq\ell\leq L.
\tag{3}
\]
For $\mathcal D\in\Delta(\mathcal X)$ and $h^\star\in\mathcal H$, draw
$x^{(0)},\ldots,x^{(T-1)}$ independently from $\mathcal D$, independently
of $\theta^{(0)}$.  With $\ell(a)=\log(1+e^{-a})$, use the exact one-sample,
all-layer update
\[
\theta^{(t+1)}=\theta^{(t)}-
\eta\nabla_\theta^{(\kappa)}
\ell\!\left(h^\star(x^{(t)})f_{\theta^{(t)}}(x^{(t)})\right),
\qquad 0\leq t<T.
\tag{4}
\]
The prescribed latter-half score, predictor, and strict classification error
are
\[
A_{\mathcal D,h^\star}(x)
:=\sum_{t=\lceil T/2\rceil}^{T}f_{\theta^{(t)}}(x),
\qquad
\widehat h_{\mathcal D,h^\star}(x)
:=\operatorname{sign}_{s_0}(A_{\mathcal D,h^\star}(x)),
\tag{5}
\]
\[
\mathcal L_{\mathcal D,h^\star}
(\widehat h_{\mathcal D,h^\star})
:=\Pr_{x\sim\mathcal D}\!\left[
\widehat h_{\mathcal D,h^\star}(x)h^\star(x)<0
\right].
\tag{6}
\]

For $r\geq0$, let
\[
B_\infty(\theta^{(0)},r)
:=\{\theta\in\mathbb R^S:
\|\theta-\theta^{(0)}\|_\infty\leq r\}.
\tag{7}
\]
Define the static hidden-preactivation margin
\[
M_r(\theta^{(0)}):=
\begin{cases}
+\infty,&L=1,\\[1mm]
\displaystyle
\inf_{\substack{\theta\in B_\infty(\theta^{(0)},r),\ x\in\mathcal X,\\
1\leq\ell<L,\ j\in[n_\ell]}}
|u_{\ell,j}(\theta,x)|,&L\geq2,
\end{cases}
\tag{8}
\]
and the exact selected-gradient envelope
\[
G_r(\theta^{(0)}):=
\sup_{\substack{\theta\in B_\infty(\theta^{(0)},r),\ x\in\mathcal X,\\
y\in\{-1,+1\}}}
\left\|\nabla_\theta^{(\kappa)}\ell(yf_\theta(x))\right\|_\infty.
\tag{9}
\]
The initialization event used below is
\[
E_r:=\{M_r(\theta^{(0)})>0,
\ \eta T G_r(\theta^{(0)})\leq r\}.
\tag{10}
\]
It is determined solely by initialization and worst-case quantities on the
fixed parameter ball, before a sample path, distribution, or target is
selected.

Finally, introduce the path index set and its dimension
\[
\mathcal I_{\rm path}:=\prod_{\ell=0}^{L-1}[n_\ell],
\qquad
d_{\rm path}:=|\mathcal I_{\rm path}|
=\prod_{\ell=0}^{L-1}n_\ell.
\tag{11}
\]
For $p=(i_0,\ldots,i_{L-1})\in\mathcal I_{\rm path}$, define the
initialization-gate feature map by
\[
[\varphi_{\theta^{(0)}}(x)]_p
:=x_{i_0}\prod_{\ell=1}^{L-1}
\mathbf 1\{u_{\ell,i_\ell}(\theta^{(0)},x)>0\},
\tag{12}
\]
where the empty product is one for $L=1$.  Let
$\mathcal P_{\rm gate}$ be the pushforward law of
$\varphi_{\theta^{(0)}}$ under the Gaussian initialization (3).  This law is
fixed by the architecture and initialization distribution before
$\mathcal D$ and $h$, and it is defined on both $E_r$ and $E_r^c$.

\subsection{Assumptions}

\begin{assumption}[Fixed source witnesses and regime]
\label{assump:fixed-source-witnesses}
The parameters satisfy $0\leq\varepsilon<1/4$, $T\in\mathbb N_{\geq1}$,
and $\eta>0$.  For the given $(n,\mathcal H,\varepsilon)$, one architecture,
its parameter count $S$, the constant stepsize $\eta$, and the horizon $T$
are selected once and held fixed before $\mathcal D$ and $h^\star$.  The same
witnesses and the same fixed $(s_0,\kappa)$ conventions are used for every
$(\mathcal D,h^\star)\in\Delta(\mathcal X)\times\mathcal H$.
\end{assumption}

\begin{assumption}[Universal expected-error SGD premise]
\label{assump:universal-expected-success}
For every $\mathcal D\in\Delta(\mathcal X)$ and every
$h^\star\in\mathcal H$,
\[
\mathbb E_{\theta^{(0)},x^{(0)},\ldots,x^{(T-1)}}\!\left[
\mathcal L_{\mathcal D,h^\star}
(\widehat h_{\mathcal D,h^\star})
\right]\leq\varepsilon,
\tag{13}
\]
where the expectation is jointly over the independent Gaussian initialization
and fresh one-sample SGD draws specified above.
\end{assumption}

\begin{assumption}[Fixed constant depth]
\label{assump:constant-depth}
A depth cap $L_0\in\mathbb N_{\geq1}$ is fixed as a universal constant,
independently of $n,\mathcal H,\varepsilon,\eta,T,S$, and the source
architecture satisfies $1\leq L\leq L_0$.  Widths remain arbitrary and all
layers remain trainable.
\end{assumption}

\begin{assumption}[Static robust initialization tube]
\label{assump:robust-tube}
Deterministic values $r>0$ and $\delta_0$ with
$0\leq\delta_0\leq\varepsilon$ are fixed after the source witnesses but
before $\theta^{(0)},\mathcal D,h^\star$, and the SGD samples, and
\[
\Pr_{\theta^{(0)}}(E_r)\geq1-\delta_0.
\tag{14}
\]
This is only a primitive statement about the static quantities $M_r$ and
$G_r$; it does not assume trajectory containment, gate stability, or a path
representation.
\end{assumption}

\subsection{Formalized goal}

The objective is the following conditional, fixed-depth probabilistic
dimension statement:
\[
\operatorname{dc}_{\varepsilon+\delta_0}(\mathcal H)
\leq d_{\rm path}\leq S^{L_0}.
\tag{15}
\]
The first inequality must be witnessed by the single target- and
distribution-independent law $\mathcal P_{\rm gate}$; equivalently, for every
$\mathcal D\in\Delta(\mathcal X)$ and $h\in\mathcal H$,
\[
\mathbb E_{\varphi\sim\mathcal P_{\rm gate}}\!\left[
\inf_{w\in\mathbb R^{d_{\rm path}}}
R_{\mathcal D,h}(w,\varphi)
\right]\leq\varepsilon+\delta_0.
\tag{16}
\]
Because $\delta_0\leq\varepsilon$, the same witness should give
\[
\operatorname{dc}_{2\varepsilon}(\mathcal H)\leq S^{L_0}.
\tag{17}
\]
The statement is conditional on Assumption~\ref{assump:robust-tube} and the
constant-depth restriction.  It does not assert the corresponding result in
the absence of the robust tube, at general depth, or with a deterministic
exact feature map.
